%=============================================================================!
% 8.S  WHAT WE NOW HAVE                                                       !
%=============================================================================!
\unnumbered{What we now have}

The potential energy on which atoms move is the energy of their electrons,
settled at each arrangement of the nuclei, plus the repulsion of the
nuclei: the potential energy surface of Born and Oppenheimer. The dense eigensolver considered here has a cost that grows as the cube
of its matrix size, motivating cheaper models of the energy. Treating the nuclei
themselves as classical is good when the quantum step $\hbar\omega$ of
their vibrations is small beside $\kB T$, which holds for heavy atoms and
slow motions and fails for the stretches of bonds to hydrogen.

The simplest models add up pair potentials, the Lennard-Jones, Morse and
exp-6 shapes, whose forces are computed once per pair and applied with
opposite signs, with the virial $\mathcal{W} = -\sum r_{ij}\varphi'(r_{ij})$
alongside; in reduced units every Lennard-Jones substance obeys the same
equations. Every force routine is checked against central differences of
its energy, at a step near the cube root of the machine epsilon times the
length scale.

A cutoff must take both energy and force smoothly to zero:
truncation breaks the conservation of energy, shifting leaves a jump in
the force that an adaptive solver follows only with far shorter steps,
and switching or a smooth envelope removes both jumps.

In metals and carbon the strength of a bond
falls as an atom's bonds multiply, which the second-moment model, the
embedded-atom method and bond-order potentials describe; molecules are
described by force fields of bonds, angles, dihedral twists and
non-bonded pairs; and the Coulomb energy of charges falls so slowly that
its sum depends on its order.

Chapter~9 uses these forces to follow the motion step by step, and
Chapter~10 builds a complete simulation of atoms in a periodic box, with
neighbour lists for the cutoffs and the Ewald sum for the charges.
